\documentclass[a4paper,landscape,headrule,footrule]{foils}
%%
%%% macros for Theories of Grammar
%%%
\usepackage{polyglossia}
\setdefaultlanguage{english}
%\setmainfont{TeX Gyre Pagella}


\newcommand{\logo}{~}
\newcommand{\header}[3]{%
\title{\vspace*{-2ex} \large JMORF/MORS --- Morpho-Syntax
\\[2ex] \Large  \red{#2} \\ \red{#3}}
\author{\blu{Francis Bond}   \\ 
\normalsize  \textbf{Palacký University}\\
\normalsize  \url{https://fcbond.github.io/}\\
\normalsize  \texttt{bond@ieee.org}}
\MyLogo{JMORF (2026)}
\renewcommand{\logo}{#2}
\hypersetup{
   pdfinfo={
     Author={Francis Bond},
     Title={#1: #2},
     Subject={JMORF/MORS: Morpho-Syntax},
     Keywords={Syntax, Morphology, HPSG, Unification, Constructions},
     License={CC BY 4.0}
   }
 }

\date{#1
  \\ Location: KB-1.11}
}
\usepackage{hyperref}
\hypersetup{
     colorlinks,
     linkcolor={blue!50!black},
     citecolor={red!50!black},
     urlcolor={blue!80!black}
   }



\usepackage{xcolor}
\usepackage{graphicx}
\newcommand{\blu}[1]{\textcolor{blue}{#1}}
\newcommand{\red}[1]{\textcolor{red}{#1}}
\newcommand{\grn}[1]{\textcolor{green}{#1}}
\newcommand{\hide}[1]{\textcolor{white}{#1}}
\newcommand{\emp}[1]{\textcolor{blue}{#1}}
\newcommand{\txx}[1]{\textbf{\textcolor{blue}{#1}}}
\newcommand{\lex}[1]{\textbf{\mtcitestyle{#1}}}


% \renewcommand{\labelitemi}{\textcolor{violet}{‣}}
% \renewcommand{\labelitemii}{\textcolor{purple}{‣}}

\newcommand{\subhead}[1]{\noindent\textbf{#1}\\[5mm]}

\newcommand{\Bad}{\emp{\raisebox{0.15ex}{\ensuremath{\mathbf{\otimes}}}}}
\newcommand{\bad}[1]{*\eng{#1}}

\newcommand{\com}[1]{\hfill (\emp{#1})}%

\usepackage{relsize,xspace}
\newcommand{\into}{\ensuremath{\rightarrow}\xspace}
\newcommand{\tot}{\ensuremath{\leftrightarrow}\xspace}
\usepackage{url}
\newcommand{\lurl}[1]{\MyLogo{\url{#1}}}

\usepackage{langsci-gb4e}
\let\eachwordone=\slshape
%%% langsci-gb4e defines \thexnumi as \@roman\@xsi{xnumi}, which only works
%%% inside an exe environment; some decks use xlisti on its own.  Restore
%%% the plain gb4e definition so that keeps working.
\makeatletter
\def\thexnumi{\@xsi{xnumi}}
\makeatother
\newcommand{\lx}[1]{\textbf{\mtciteform{#1}}}
\newcommand{\ix}{\ex\slshape}



\newcommand{\ent}{\ensuremath{\Rightarrow}\xspace}
\newcommand{\ngv}{\ensuremath{\not\Rightarrow}\xspace}
%\usepackage{times}
%\usepackage{nttfoilhead}
\newcommand{\myslide}[1]{\foilhead[-25mm]{\raisebox{12mm}[0mm]{\red{#1}}}\MyLogo{\logo}}
\newcommand{\myslider}[1]{\rotatefoilhead[-25mm]{\raisebox{12mm}[0mm]{\red{#1}}}}
%\newcommand{\myslider}[1]{\rotatefoilhead{\raisebox{-8mm}{\emp{#1}}}}

\newcommand{\section}[1]{\myslide{}{\begin{center}\Huge \red{#1}\end{center}}}



\usepackage[lyons,j,e,k]{mtg2e}
%\renewcommand{\mtcitestyle}[1]{\textcolor{-red!75!green!50}{\textsl{#1}}}
\renewcommand{\mtcitestyle}[1]{\textcolor{teal}{\textsl{#1}}}
\newcommand{\cs}{\mtciteform}
\newcommand{\iz}[1]{\texttt{\textup{#1}}}
\newcommand{\gm}{\textsc}
\usepackage[normalem]{ulem}
\newcommand{\ul}{\uline}
\newcommand{\ull}{\uuline}
\newcommand{\wl}{\uwave}
\newcommand{\vs}{\ensuremath{\Leftrightarrow}~}
%%%
%%% Bibliography
%%%
\usepackage{natbib}
%\usepackage{url}
\usepackage{bibentry}



\usepackage[utf8]{inputenc}
%%% trees
\usepackage{multicol}
%%% pst-node: one example in lec-12-control draws a connector between
%%% two points with \rnode/\ncbar.  This used to come in with rtrees.
\usepackage{pst-node}
\usepackage{forest}
\useforestlibrary{linguistics}
\forestapplylibrarydefaults{linguistics}

%%% \makebox's optional arguments are read as AVM delimiters by
%%% langsci-avm, so hide them behind a macro with no visible brackets.
\newcommand{\lbox}[1]{\makebox[0mm][l]{#1}}

%%% Tree styles.
%%%   words  -- trees whose leaves are actual words (italicised, aligned)
%%%   cat    -- trees whose leaves are categories or feature structures
\forestset{
  words/.style={for tree={s sep=7mm, l sep=5mm,
                          if n children=0{font=\itshape, tier=terminal}{}}},
  cat/.style={for tree={s sep=7mm, l sep=5mm}},
}
\newcommand{\wtree}[1]{\begin{forest} words #1 \end{forest}}
\newcommand{\ctree}[1]{\begin{forest} cat #1 \end{forest}}
%%% the rule being applied at this step of a derivation
\newcommand{\step}[1]{\begin{center}\emp{#1}\end{center}}

%%% AVMs (Language Science Press)
\usepackage{langsci-avm}

%%% Macros the decks use, carried over from headx.tex
\newcommand{\blank}{\rule{3em}{1pt}\xspace}
\newcommand{\ft}[1]{\textsc{#1}}
\newcommand{\idn}{\mtciteform}    %%% Indonesian
\newcommand{\prd}[1]{\textbf{#1}}
\newcommand{\typ}[1]{\textit{#1}}
%%% \val: provide before renewing, in case a package defines it.
\providecommand{\val}[1]{}
\renewcommand{\val}[1]{\textit{#1}}
%%% Manning's \avmbox, still used in tree labels and in maths where a
%%% bare langsci \tag would not work.  \avm{\tag{...}} gives the same
%%% boxed tag as the ones inside feature structures.
\newcommand{\avmbox}[1]{\mbox{\avm{\tag{#1}}}}

%%% Ring a value in a dotted oval.  These used pstricks \ovalnode,
%%% which came in with rtrees; tikz (via forest) does the same job.
\newcommand{\OV}[1]{\tikz[baseline]{\node[draw,dotted,red,ellipse,
          inner sep=1pt,outer sep=0pt]{#1};}}
\newcommand{\OVB}[1]{\tikz[baseline]{\node[draw,dotted,blue,ellipse,
          inner sep=1pt,outer sep=0pt]{#1};}}

%%% From CSLI book
\newcommand{\mc}{\multicolumn}
\newcommand{\HD}{\textbf{H}\xspace}
\newcommand{\el}{< >}       %%% empty list; langsci-avm needs the space

\newcommand{\A}{\noindent\textbf{A}: }
\newcommand{\Q}{\noindent\textbf{Q}: }

%%% Question and answer on two pages, written once:
%%%   \qaslide{Title}{... question ... \answer{... answer ...}}
%%% prints the slide twice, first without the \answer parts and then
%%% with them, so the answer appears "on click".  The example counter
%%% is reset so both pages show the same example numbers: gb4e keeps the
%%% running number in the equation counter between exe environments, and
%%% in xnumi inside one.
\newif\ifshowanswer
\newcommand{\answer}[1]{\ifshowanswer#1\fi}
\newcounter{qaxnumi}
\newcounter{qaeqn}
\newcommand{\qaslide}[2]{%
  \setcounter{qaxnumi}{\value{xnumi}}%
  \setcounter{qaeqn}{\value{equation}}%
  \showanswerfalse\myslide{#1}#2%
  \setcounter{xnumi}{\value{qaxnumi}}%
  \setcounter{equation}{\value{qaeqn}}%
  \showanswertrue\myslide{#1}#2%
  \showanswerfalse}
%\newcommand{\C}{\noindent\textbf{C}: }



%%% Extra lecture: students may read this on their own, so every
%%% question slide carries its answer (\qaslide ... \answer).
\begin{document}
\header{Lecture 13}{Auxiliaries}{}
\maketitle

\overview
\MyLogo{\citet{sag:wasow:bender:2003} --- Chapter 13}

\section[Introduction: what auxiliaries are and how they behave]{Introduction}

\qaslide{Where We Are \& Where We’re Going}{
\begin{itemize}
\item Revision
  \begin{itemize}
  \item Last time we met \txx{subject-raising verbs} (\val{srv-lxm}).
    What is special about their subject?
    \answer{\\ {\small\emp{It is also the subject of the complement VP
      (tag \avmbox{1}), and the verb gives it no semantic role of its
      own: \eng{There continues to be trouble}.}}}
  \item We also called infinitival \lex{to} an \txx{auxiliary verb},
    [\ft{aux} $+$], and promised the reason would come later.
  \end{itemize}
\item Today: the real auxiliaries, \lex{be}, \lex{have}, \lex{do} and
  the modals (\lex{can}, \lex{will}, \lex{must}, \ldots).
\item They turn out to be subject-raising verbs with a few extra
  features and four lexical rules.
\item This is an extra lecture: read it at your own pace, and try each
  question before turning the page.
\end{itemize}
}

\myslide{What Auxiliaries Are}
\begin{itemize}
\item Sometimes called ``helping verbs'', auxiliaries are little words
  that come before the main verb of a sentence, including forms of
  \lex{be}, \lex{have}, \lex{do}, \lex{can}, \lex{could}, \lex{may},
  \lex{might}, \lex{must}, \lex{shall}, \lex{should}, \lex{will} and
  \lex{would}.
\item They tend to be involved in the expression of time, necessity,
  possibility, permission and obligation, as well as such things as
  negation, affirmation and questioning.
\end{itemize}

\myslide{Some Basic Facts about Auxiliaries}
\begin{itemize}\addtolength{\itemsep}{-0.5ex}
\item They are optional \\
  \eng{Pat tapdanced.} \eng{Pat can tapdance.} \eng{Pat is tapdancing.}
\item They precede any non-auxiliary verbs \\
  \bad{Pat tapdance can.} \bad{Pat tapdancing is.}
\item They determine the form of the following verb \\
  \bad{Pat can tapdancing.} \bad{Pat is tapdance.}
\item When they co-occur, their order is fixed \\
  \eng{Pat must be tapdancing.} \bad{Pat is musting tapdance.}
\item Auxiliaries of any given type cannot iterate \\
  \bad{Pat could should tapdance.}
\end{itemize}

\myslide{A Little History}
\begin{itemize}
\item Chomsky’s first book, \textit{Syntactic Structures} (1957),
  contained a detailed analysis of the English system of auxiliary
  verbs.
\item It showed how formal analysis could reveal subtle
  generalizations.
\item The power of Chomsky’s analysis of auxiliaries was one of the
  early selling points for transformational grammar
  \begin{itemize}
  \item especially his unified treatment of auxiliary \lex{do}.
  \end{itemize}
\item So it’s a challenge to any theory of grammar to deal with the
  same phenomena.
\end{itemize}

\myslide{Two Approaches to Analyzing Auxiliaries}
\begin{itemize}
\item Treat auxiliaries as a special category, and formulate
  specialized transformations sensitive to their presence.
\item Assimilate their properties to existing types as much as
  possible, and elaborate the lexicon to handle what is special about
  them.
\item We adopt the latter, treating auxiliaries as a subtype of
  \val{srv-lxm}.
\end{itemize}

\qaslide{Consequences of Making \typ{auxv-lxm} a Subtype of \typ{srv-lxm}}{
\begin{itemize}
\item Auxiliaries should express one-place predicates.
\item Auxiliaries should allow non-referential subjects (dummy
  \lex{there}, \lex{it}, and idiom chunks).
\item Passivization of the main verb (the auxiliary’s complement)
  should preserve truth conditions.
\item Are these borne out?
  \answer{\\ {\small\emp{Yes:
    \eng{There might be a unicorn in the garden},
    \eng{It will annoy people that dogs are barking},
    \eng{Close tabs have been kept on Pat by the FBI};
    and \eng{A doctor must have been examining Pat} describes the same
    situations as \eng{Pat must have been being examined by a doctor}.}}}
\end{itemize}
}

\myslide{Why Call Auxiliaries Verbs?}
\begin{itemize}
\item \lex{be}, \lex{have} and \lex{do} exhibit verbal inflections
  (tense, agreement).
\item \lex{be}, \lex{have} and \lex{do} can all appear as main verbs
  (that is, as the only verb in a sentence).
\item Their inflections are the same in main and auxiliary uses.
\item \lex{be} exhibits auxiliary behaviour even in its main verb
  uses: \eng{Kim is not tall}, \eng{Is Kim tall?}
\item Modals (\lex{can}, \lex{might}, \lex{will}, \ldots) don’t
  inflect, but they occur in environments requiring a finite verb with
  no (other) finite verb around.
\end{itemize}

\myslide{What’s Special about Auxiliaries?}
\begin{itemize}
\item Unlike other subject-raising verbs we have looked at, their
  complements aren’t introduced by \lex{to}.
\item The modals and \lex{do} have defective paradigms (no
  \bad{to can}, \bad{musting}).
\item There are restrictions on the ordering and iterability of
  auxiliaries.
\item They have a set of special characteristics known as the
  \txx{NICE} properties (Negation, Inversion, Contraction, Ellipsis).
\end{itemize}

\section[Lexical types and entries: \typ{auxv-lxm}, \lex{be}, \lex{have}, the modals]{Lexical Types and Entries}

\myslide{Where \typ{auxv-lxm} Sits}
\begin{center}
\begin{forest}
cat, for tree={font=\itshape, l sep=6mm, s sep=4mm}
[{verb-lxm}
  [{srv-lxm} [{ic-srv-lxm}] [{auxv-lxm}]]
  [{scv-lxm}] [{piv-lxm}] [{siv-lxm}]
  [{tv-lxm} [{\ldots}]]
  [{dervv-lxm}]]
\end{forest}
\end{center}
\begin{itemize}\addtolength{\itemsep}{-1ex}
\item \val{srv-lxm} splits in two: \val{ic-srv-lxm} is the
  \lex{continue} type from last week (infinitival complement),
  \val{auxv-lxm} is new.
\item \val{dervv-lxm} is for the outputs of the Ellipsis LR: we get
  to it at the end.
\end{itemize}

\myslide{Some Type Constraints (1)}
\begin{center}\small
\begin{tabular}{|l|l|l|}
\hline
\textsc{type} & \textsc{features/constraints} & \textsc{ist} \\ \hline
\val{verb-lxm} &
  \avm{[ syn & [ head & [\type*{verb} aux & / $-$ ] ] \\
         arg-st & < [ head & nominal ], \ldots > \\
         sem & [ mode & prop ] ]} & \val{infl-lxm} \\ \hline
\val{srv-lxm} &
  \avm{[ arg-st & < \1, [ spr & < \1 > \\ comps & < > ] > ]} &
  \val{verb-lxm} \\ \hline
\end{tabular}
\end{center}
\begin{itemize}\addtolength{\itemsep}{-1ex}
\item Nothing new here except the default [\ft{aux} $/-$] on every verb.
\end{itemize}

\myslide{Some Type Constraints (2)}
\begin{center}\small
\begin{tabular}{|l|l|l|}
\hline
\textsc{type} & \textsc{features/constraints} & \textsc{ist} \\ \hline
\val{ic-srv-lxm} &
  \avm{[ arg-st & < X, VP [ inf & $+$ \\ index & s ] > \\
         sem & [ restr & < [ arg & s ] > ] ]} & \val{srv-lxm} \\ \hline
\val{auxv-lxm} &
  \avm{[ syn & [ head & [ aux & $+$ ] ] ]} & \val{srv-lxm} \\ \hline
\end{tabular}
\end{center}
\begin{itemize}\addtolength{\itemsep}{-1ex}
\item \val{ic-srv-lxm}: last week's \lex{continue}, \lex{tend},
  \lex{seem} (``infinitival-complement'').
\item \val{auxv-lxm} adds just one thing: it overrides the default and
  makes the verb [\ft{aux} $+$].  Everything else comes from
  \val{srv-lxm}.
\end{itemize}

\myslide{A Lexical Entry for \lex{be}}
\begin{center}
  \avm{< \textnormal{be}, [\type*{auxv-lxm}
    arg-st & < X, [ syn & [ head & [ pred & $+$ ] ] \\
                    sem & [ index & \2 ] ] > \\
    sem & [ index & \2 \\ restr & < > ] ] >}
\end{center}
\begin{itemize}\addtolength{\itemsep}{-1ex}
\item The same \lex{be} as in Chapter 11: a [\ft{pred} $+$]
  complement, no meaning of its own.
\item Now it has a proper type.  Compare the entry with inherited
  information on the next slide.
\end{itemize}

\qaslide{The Entry for \lex{be}, with Inherited Information}{
\begin{center}\footnotesize
  \avm{< \textnormal{be}, [\type*{auxv-lxm}
    syn & [ head & [\type*{verb} aux & $+$ \\ agr & \tag{0} ] \\
            val & [ spr & < [ agr & \tag{0} ] > ] ] \\
    arg-st & < \3, [ syn & [ head & [ pred & $+$ ] \\
                              val & [ spr & < \3 > \\ comps & < > ] ] \\
                     sem & [ index & \2 ] ] > \\
    sem & [ mode & prop \\ index & \2 \\ restr & < > ] ] >}
\end{center}
\begin{itemize}\addtolength{\itemsep}{-1ex}
\item Where does each piece come from?
  \answer{\\ {\small\emp{[\ft{aux} $+$]: \val{auxv-lxm}.  The shared
    subject \avmbox{3} and empty \ft{comps}: \val{srv-lxm}.
    [\ft{mode} prop] and the \ft{agr} sharing with the \ft{spr}:
    \val{verb-lxm} and the SHAC.  Only [\ft{pred} $+$] and the
    \ft{index} sharing are \lex{be}'s own.}}}
\end{itemize}
}

\qaslide{Entry for \lex{have}}{
\begin{center}\small
  \avm{< \textnormal{have}, [\type*{auxv-lxm}
    arg-st & < X, [ syn & [ head & [\type*{verb} form & psp ] ] \\
                    sem & [ index & \3 ] ] > \\
    sem & [ index & s \\
            restr & < [ reln & \textbf{have} \\ sit & s \\ arg & \3 ] > ] ] >}
\end{center}
\begin{itemize}\addtolength{\itemsep}{-1ex}
\item Note the \ft{form} restriction on the complement VP: a past
  participle, \eng{has \ul{eaten}}.
\item What accounts for the analogous \ft{form} restriction on verbs
  following \lex{be}?
  \answer{\\ {\small\emp{\lex{be} asks for [\ft{pred} $+$].  Verbs are
    [\ft{pred} $-$] by default; only the Present Participle LR and the
    Passive LR make a verb [\ft{pred} $+$].  So after \lex{be} we get
    \eng{eating} or \eng{eaten} (passive) but not \bad{is eat}.}}}
\end{itemize}
}

\qaslide{Lexical Entry for a Modal}{
\begin{center}\footnotesize
  \avm{< \textnormal{would}, [\type*{auxv-lxm}
    syn & [ head & [ form & fin ] ] \\
    arg-st & < X, [ syn & [ head & [\type*{verb} inf & $-$ \\ form & base ] ] \\
                    sem & [ index & s$_2$ ] ] > \\
    sem & [ index & s$_1$ \\
            restr & < [ reln & \textbf{would} \\ sit & s$_1$ \\ arg & s$_2$ ] > ] ] >}
\end{center}
\begin{itemize}\addtolength{\itemsep}{-1ex}
\item Note the restriction on the form of the complement VP: base
  form, not \lex{to} ([\ft{inf} $-$]).
\item What inflectional lexical rules apply to this lexeme?
  \answer{\\ {\small\emp{Only ones that keep [\ft{form} fin]: the
    3rd-Singular and Non-3rd-Singular Verb LRs (F$_{3SG}$ leaves modals
    alone: \lex{would}, not \bad{woulds}).  Not the participle or Base
    Form LRs, whose outputs have another \ft{form} value.}}}
\end{itemize}
}

\myslide{Accounting for the Basic Facts Cited Earlier}
\begin{itemize}\addtolength{\itemsep}{-0.5ex}
\item \txx{Optionality of auxiliaries:} as raising verbs, their
  subjects and complements go together.
\item \txx{Auxiliaries precede non-auxiliary verbs:} auxiliaries are
  heads, and complements follow heads in English.
\item \txx{Auxiliaries determine the form of the following verb:}
  this is built into their lexical entries.
\item \txx{When auxiliaries co-occur, their order is fixed:}
  different explanations for different combinations; see next slide.
\item \txx{Non-iterability of auxiliaries:} ditto.
\end{itemize}

\myslide{Accounting for Restrictions on Order and Iterability}
\begin{itemize}
\item Order
  \begin{itemize}
  \item Modals are finite, and all auxiliaries take non-finite
    complements.  Hence modals must come first.
  \item Stative verbs (like \lex{own}) don’t have present participles,
    and auxiliary \lex{have} is stative.  Hence \bad{Pat is having
      tapdanced}.
  \end{itemize}
\item Iterability
  \begin{itemize}
  \item Auxiliary \lex{be} is also stative, so \bad{Pat is being
      tapdancing}.
  \item Modals must be finite, and their complements must be base, so
    \bad{Pat can should tapdance}.
  \item \bad{Pat has had tapdanced} can be ruled out in various ways,
    e.g.\ stipulating that auxiliary \lex{have} has no past participle.
  \end{itemize}
\end{itemize}

\section[The old transformational analysis, and a constituency argument against it]{Comparison: The Old Transformational Analysis}

\myslide{Sketch of Chomsky’s Old Analysis}
\begin{center}
S \into\ NP AUX VP \qquad AUX \into\ T (M) (PERF) (PROG)

\medskip
\begin{forest}
cat, for tree={s sep=5mm}
[{S} [{NP} [{Chris}]]
  [{AUX} [{T} [{\red{past}}]] [{M} [{could}]] [{PERF} [{have+\red{en}}]] [{PROG} [{be+\red{ing}}]]]
  [{VP} [{V} [{eat}]]]]
\end{forest}
\end{center}
\begin{itemize}\addtolength{\itemsep}{-1ex}
\item A flat AUX node holds tense and the auxiliaries; each red
  affix then hops onto the verb to its right: \eng{Chris could have
    been eating}.
\end{itemize}

\myslide{How This Analysis Handles the Basic Facts}
\begin{itemize}\addtolength{\itemsep}{-0.5ex}
\item \txx{Optionality of auxiliaries:} stipulated in the phrase
  structure rule (with parentheses).
\item \txx{Auxiliaries precede non-auxiliary verbs:} built into the
  phrase structure rule, with AUX before VP.
\item \txx{Auxiliaries determine the form of the following verb:}
  inflections are inserted with the auxiliaries and moved onto the
  following verb transformationally.
\item \txx{When auxiliaries co-occur, their order is fixed:}
  stipulated in the phrase structure rule for AUX.
\item \txx{Non-iterability of auxiliaries:} ditto.
\end{itemize}

\myslide{The Two Analyses Assign Very Different Trees}
\begin{minipage}[t]{0.48\textwidth}\vspace{0pt}\centering
\scalebox{0.8}{\begin{forest}
cat
[{S} [{NP}] [{VP} [{V} [{could}]] [{VP} [{V} [{have}]] [{VP} [{V} [{been}]] [{VP}]]]]]
\end{forest}}

\small
\begin{itemize}\addtolength{\itemsep}{-1ex}
\item \eng{could have been} VP, \eng{have been} VP and \eng{been} VP
  are all constituents
\item \eng{could have been} is not a constituent
\end{itemize}
\end{minipage}\hfill
\begin{minipage}[t]{0.48\textwidth}\vspace{0pt}\centering
\scalebox{0.8}{\begin{forest}
cat
[{S} [{NP}] [{AUX} [{M} [{could}]] [{PERF} [{have}]] [{PROG} [{been}]]] [{VP}]]
\end{forest}}

\small
\begin{itemize}\addtolength{\itemsep}{-1ex}
\item \eng{could have been} VP, \eng{have been} VP and \eng{been} VP
  are not constituents
\item \eng{could have been} is a constituent
\end{itemize}
\end{minipage}

\qaslide{Ellipsis and Constituency}{
\begin{itemize}
\item Consider:
  \begin{exe}
  \ex \eng{Pat couldn’t have been eating garlic, but Chris could have been.}
  \ex \eng{Pat couldn’t have been eating garlic, but Chris could have.}
  \ex \eng{Pat couldn’t have been eating garlic, but Chris could.}
  \end{exe}
\item What is missing in each case, and which analysis makes it a
  constituent?
  \answer{\\ {\small\emp{\eng{eating garlic}, \eng{been eating garlic},
    \eng{have been eating garlic}.  On the nested analysis each is a
    VP; on the flat analysis none of them is a constituent.  Ellipsis
    normally removes constituents, so this argues for our analysis over
    the old transformational one.}}}
\end{itemize}
}

\myslide{Our Analysis of Auxiliaries So Far}
\begin{itemize}
\item Auxiliaries are subject-raising verbs.
\item Most basic distributional facts about them can be handled
  through selectional restrictions between auxiliaries and their
  complements (that is, as \ft{arg-st} constraints).
\item Auxiliaries are identified via a \ft{head} feature \ft{aux},
  which we have not yet put to use.
\end{itemize}

\section[The NICE properties: negation, inversion, contraction and ellipsis as four lexical rules]{The NICE Properties}

\myslide{Descriptive Summary of the NICE Properties}
\begin{center}
\begin{tabular}{lp{0.7\textwidth}}
\txx{Negation} & Sentences are negated by putting \lex{not} after the
  first auxiliary verb; they can be reaffirmed by putting \lex{too} or
  \lex{so} in the same position. \\[1ex]
\txx{Inversion} & Questions are formed by putting an auxiliary verb
  before the subject NP. \\[1ex]
\txx{Contraction} & Auxiliary verbs take negated forms, with
  \lex{n’t} affixed. \\[1ex]
\txx{Ellipsis} & Verb phrases immediately following an auxiliary verb
  can be omitted. \\
\end{tabular}
\end{center}

\qaslide{Negation (and Reaffirmation)}{
\begin{itemize}
\item Polar adverbs (sentential \lex{not}, \lex{so} and \lex{too})
  appear immediately following an auxiliary:
  \begin{exe}
  \ex \eng{Pat will not leave.} \quad \eng{Pat will SO leave.} \quad
      \eng{Pat will TOO leave.}
  \end{exe}
\item What about examples like \eng{Not many people left}?
  \answer{\\ {\small\emp{That is constituent negation: \lex{not}
    modifies the NP \eng{many people}, not the sentence.  We set it
    aside.}}}
\item What happens when you want to deny or reaffirm a sentence with
  no auxiliary?
  \begin{exe}
  \ex \eng{Pat left.} \quad \eng{Pat did not leave.} \quad
      \eng{Pat did TOO leave.}
  \end{exe}
  \answer{{\small\emp{An auxiliary appears from nowhere: \lex{do}.}}}
\end{itemize}
}

\myslide{The Auxiliary \lex{do}}
\begin{itemize}\addtolength{\itemsep}{-1ex}
\item Like modals, auxiliary \lex{do} only occurs in finite contexts:
  \bad{Pat continued to do not leave}.
\item Unlike modals, \lex{do} cannot be followed by other auxiliaries:
  \bad{Pat did not have left}.
\end{itemize}
\begin{center}\small
  \avm{< \textnormal{do}, [\type*{auxv-lxm}
    syn & [ head & [ form & fin ] ] \\
    arg-st & < X, [ syn & [ head & [\type*{verb} form & base \\ aux & $-$ ] ] \\
                    sem & [ index & s ] ] > \\
    sem & [ index & s \\ restr & < > ] ] >}
\end{center}
\emp{[\ft{aux} $-$] on the complement does the second job; like \lex{be},
  \lex{do} adds no meaning of its own.}

\myslide{The ADV$_{pol}$-Addition Lexical Rule}
\begin{center}\small
  \avm{[\type*{pi-rule}
    input & < X, [ syn & [ head & [\type*{verb} form & fin \\ pol & $-$ \\ aux & $+$ ] ] \\
                   arg-st & < \1 > $\oplus$ \tag{A} \\
                   sem & [ index & s$_1$ ] ] > \\
    output & < Y, [ syn & [ head & [ pol & $+$ ] \\ val & [ spr & < Z > ] ] \\
                    arg-st & < \1 > $\oplus$ < ADV$_{pol}$ [ index & s$_2$ \\
                                                             restr & < [ arg & s$_1$ ] > ] > $\oplus$ \tag{A} \\
                    sem & [ index & s$_2$ ] ] > ]}
\end{center}
\emp{Take a finite auxiliary word and add a polar adverb right after
  its subject on \ft{arg-st}; the adverb's predication takes the
  auxiliary's situation as its argument.}

\myslide{What Does the Type \typ{pi-rule} Mean?}
\begin{itemize}\addtolength{\itemsep}{-1ex}
\item It maps words to words (hence ``post-inflectional'').
\item It preserves \ft{mod} values, \ft{head} values (defeasibly) and
  \ft{sem} values (defeasibly).
\end{itemize}
\begin{center}
  \avm{[ input & < / \tag{0}, [\type*{word}
                   syn & [ head & / \1 \\ val & [ mod & \tag{A} ] ] \\
                   sem & / \2 ] > \\
         output & < / \tag{0}, [\type*{word}
                   syn & [ head & / \1 \\ val & [ mod & \tag{A} ] ] \\
                   sem & / \2 ] > ]}
\end{center}
\emp{So \ft{form} fin and \ft{aux} $+$ carry over to the output without
  being mentioned; only \ft{pol} and \ft{index} are overridden.}

\qaslide{Why Doesn’t the ADV$_{pol}$-Addition LR Mention \ft{comps}?}{
\begin{center}\scriptsize
  \avm{[\type*{pi-rule}
    input & < X, [ syn & [ head & [\type*{verb} form & fin \\ pol & $-$ \\ aux & $+$ ] ] \\
                   arg-st & < \1 > $\oplus$ \tag{A} \\
                   sem & [ index & s$_1$ ] ] > \\
    output & < Y, [ syn & [ head & [ pol & $+$ ] \\ val & [ spr & < Z > ] ] \\
                    arg-st & < \1 > $\oplus$ < ADV$_{pol}$ [ index & s$_2$ \\
                                                             restr & < [ arg & s$_1$ ] > ] > $\oplus$ \tag{A} \\
                    sem & [ index & s$_2$ ] ] > ]}
\end{center}
\begin{itemize}\addtolength{\itemsep}{-1ex}
\item The rule changes \ft{arg-st} but says nothing about \ft{comps}.  Why not?
  \answer{\\ {\small\emp{The output is a word, so the Argument
    Realization Principle fills in \ft{val} from \ft{arg-st}.  With
    [\ft{spr} $\langle$Z$\rangle$] the first argument \avmbox{1} is the
    specifier, and everything else, the new ADV$_{pol}$ first, lands on
    \ft{comps}.}}}
\end{itemize}
}

\qaslide{What Is the Role of These Indices?}{
\begin{center}\scriptsize
  \avm{[\type*{pi-rule}
    input & < X, [ syn & [ head & [\type*{verb} form & fin \\ pol & $-$ \\ aux & $+$ ] ] \\
                   arg-st & < \1 > $\oplus$ \tag{A} \\
                   sem & [ index & \OV{s$_1$} ] ] > \\
    output & < Y, [ syn & [ head & [ pol & $+$ ] \\ val & [ spr & < Z > ] ] \\
                    arg-st & < \1 > $\oplus$ < ADV$_{pol}$ [ index & \OVB{s$_2$} \\
                                                             restr & < [ arg & \OV{s$_1$} ] > ] > $\oplus$ \tag{A} \\
                    sem & [ index & \OVB{s$_2$} ] ] > ]}
\end{center}
\begin{itemize}\addtolength{\itemsep}{-1ex}
\item $s_1$ and $s_2$: what do they do?
  \answer{\\ {\small\emp{$s_1$ is the auxiliary's situation; the
    adverb takes it as its \ft{arg}: \textbf{not}($s_1$).  $s_2$ is the
    adverb's situation, and the output takes it as its \ft{index}, so
    the negation becomes the main predication of the sentence.}}}
\end{itemize}
}

\myslide{The Entry for \lex{not}}
\begin{center}
  \avm{< \textnormal{not}, ADV$_{pol}$ [ sem & [ index & s$_1$ \\
      restr & < [ reln & \textbf{not} \\ sit & s$_1$ \\ arg & s$_2$ ] > ] ] >}
\end{center}
\begin{itemize}\addtolength{\itemsep}{-1ex}
\item An ADV$_{pol}$ is an adverb word with a one-place predication
  and an index of its own.
\item The rule on the previous slide unifies its \ft{arg} with the
  auxiliary's index and passes its own index up.
\item \lex{so} and \lex{too} are the same, with \textbf{so} and
  \textbf{too} as their relations.
\end{itemize}

\myslide{``Scope'' of Negation with the ADV$_{pol}$-Addition LR}
\begin{itemize}\addtolength{\itemsep}{-0.5ex}
\item Note that the formulation of this rule assumes that \lex{not}
  takes ``wide scope''.
\item \eng{You can not fall asleep}
  \begin{itemize}
  \item \lex{not} takes \lex{can} as a semantic argument,
  \item resulting in the meaning where it is not the case that you are
    able to fall asleep.
  \end{itemize}
\item vs.\ \eng{You must not fall asleep}
  \begin{itemize}
  \item here, the desired interpretation is ``you are not allowed to
    fall asleep'',
  \item however, that interpretation requires \lex{must} to take scope
    over \lex{not},
  \item so to properly handle this we would need another rule, but
    this grammar fragment does not.
  \end{itemize}
\item This lexical difference is part of why we treat \lex{not} as a
  complement and not a modifier.
\end{itemize}

\qaslide{Which \lex{not}s Does the Rule License?}{
\begin{exe}
\ex \eng{Andy must \ul{not} have been sleeping.}
\ex \eng{Andy must have \ul{not} been sleeping.}
\ex \eng{Andy must have been \ul{not} sleeping.}
\ex \eng{Kleptomaniacs can\ul{not} not steal.}
\ex \eng{Kleptomaniacs cannot \ul{not} steal.}
\end{exe}
\begin{itemize}
\item For each underlined \lex{not}: is it the ADV$_{pol}$ complement
  of a finite auxiliary?
  \answer{\\ {\small\emp{(1) yes: first complement of finite
    \lex{must}.  (2), (3) no: \lex{have} and \lex{been} are not
    finite, so these are constituent negation.  (4) yes.  (5) no: the
    rule needs a [\ft{pol} $-$] input and \lex{can} has already become
    [\ft{pol} $+$], so the second \lex{not} is constituent negation of
    \eng{steal}.  Key point: sentential \lex{not} must show up as the
    first complement of a finite auxiliary.}}}
\end{itemize}
}

\myslide{Negation and Reaffirmation: Sample Tree}
\begin{center}
\begin{forest}
cat
[{S} [{NP} [{Leslie}]]
  [{VP} [{V} [{did}]] [{ADV$_{pol}$} [{SO}]] [{VP} [{eat the whole pizza}, roof]]]]
\end{forest}
\end{center}
\begin{itemize}\addtolength{\itemsep}{-1ex}
\item \lex{did} has three arguments: the subject, the ADV$_{pol}$ and
  the base VP.  The lower VP is licensed by the Head-Complement Rule
  with all of them.
\end{itemize}

\qaslide{Inversion}{
\begin{itemize}
\item Yes-no questions begin with an auxiliary: \eng{Will Robin win?}
\item The NP after the auxiliary has all the properties of a subject
  \begin{itemize}
  \item Agreement: \eng{Have they left?} vs.\ \bad{Has they left?}
  \item Case: \bad{Have them left?}
  \item Raising: \eng{Will there continue to be food at the meetings?}
  \end{itemize}
\item What happens if you make a question out of a sentence without
  an auxiliary?  \eng{Robin won.} \into\ ?
  \answer{\\ {\small\emp{\eng{Did Robin win?}  Again \lex{do} steps in,
    and \eng{Robin} still agrees with it: \eng{Does Robin win?}}}}
\end{itemize}
}

\myslide{The Inversion Lexical Rule}
\begin{center}
  \avm{[\type*{pi-rule}
    input & < W, [ syn & [ head & [\type*{verb} form & fin \\ aux & $+$ ] ] \\
                   arg-st & \tag{A} \\
                   sem & [ mode & prop ] ] > \\
    output & < Z, [ syn & [ head & [ inv & $+$ ] \\ val & [ spr & < > ] ] \\
                    arg-st & \tag{A} \\
                    sem & [ mode & ques ] ] > ]}
\end{center}
\emp{Take a finite auxiliary word, keep its arguments, empty its
  \ft{spr} list and make it a question.}

\qaslide{How the Rule Yields Inverted Order}{
\begin{center}\scriptsize
  \avm{[\type*{pi-rule}
    input & < W, [ syn & [ head & [\type*{verb} form & fin \\ aux & $+$ ] \\
                           val & [ spr & \OV{< X >} ] ] \\
                   arg-st & \tag{A} \\
                   sem & [ mode & prop ] ] > \\
    output & < Z, [ syn & [ head & [ inv & $+$ ] \\ val & [ spr & \OVB{< >} ] ] \\
                    arg-st & \tag{A} \\
                    sem & [ mode & ques ] ] > ]}
\end{center}
\begin{itemize}\addtolength{\itemsep}{-1ex}
\item The \ft{arg-st} is unchanged.  So how does the subject end up
  after the verb?
  \answer{\\ {\small\emp{The ARP again: with \ft{spr} empty, every
    argument including the subject is on \ft{comps}, and the
    Head-Complement Rule puts complements after the head: V NP VP.  The
    result has empty \ft{spr} and \ft{comps}, so it is an S.}}}
\end{itemize}
}

\qaslide{The Feature \ft{inv}}{
\begin{itemize}
\item What is the \ft{inv} value of inputs to the Inversion LR?
  \answer{\\ {\small\emp{Perhaps surprisingly, the input is also
    [\ft{inv} $+$]: \val{pi-rule}s have default identity of \ft{head}
    features, and the input gives no \ft{inv} value.}}}
\item Then what work is the feature doing?
  \answer{\\ {\small\emp{It marks auxiliaries that can't be inverted
    against those that must be:
    \eng{You better watch out} vs.\ \bad{Better you watch out}
    (\lex{better} is [\ft{inv} $-$], so the rule cannot apply);
    \eng{I shall go} (\lex{shall} $\approx$ `will') vs.\ \eng{Shall I
      go?} (\lex{shall} $\approx$ `should'): two different words, one
    [\ft{inv} $-$], one [\ft{inv} $+$].}}}
\end{itemize}
}

\qaslide{Other Cases of Inversion}{
\begin{itemize}
\item Inversion is not limited to questions:
  \begin{itemize}
  \item Preposed negatives: \eng{Never have I been so upset!}
  \item Conditionals: \eng{Had we known, we would have left.}
  \item Exclamations: \eng{May your teeth fall out!}
  \end{itemize}
\item Does our rule account for these?
  \answer{\\ {\small\emp{No.  Our rule's output says [\ft{mode} ques],
    and each construction has slightly different idiosyncrasies.}}}
\item How might we extend our analysis to cover them?
  \answer{\\ {\small\emp{One option: define a type of inversion lexical
    rules, sharing certain properties but with some differences.
    Another: define more phrase structure rules that select for
    [\ft{inv} $+$] things.}}}
\end{itemize}
}

\myslide{Inversion: Sample Tree}
\begin{center}
\begin{forest}
cat
[{S} [{V} [{Did}]] [{NP} [{Leslie}]] [{VP} [{eat the entire pizza?}, roof]]]
\end{forest}
\end{center}
\begin{itemize}\addtolength{\itemsep}{-1ex}
\item One flat S: the Head-Complement Rule with \lex{Did} and both of
  its complements, the NP and the VP.  No Head-Specifier Rule is
  involved.
\end{itemize}

\myslide{Contraction}
\begin{itemize}
\item There are several types of contraction in English, but we’re
  only talking about words ending in \lex{n’t}.
\item It may seem like it’s just \lex{not} said fast, but there’s more
  to it:
  \begin{itemize}
  \item Only finite verbs can take \lex{n’t}: \bad{Terry must haven’t
      seen us}.  If it was just \lex{not} said fast you’d expect to be
    able to say this.
  \item There are morphological irregularities: \lex{won’t} not
    \bad{willn’t}; \%\lex{shan’t}, not \bad{shalln’t}; \lex{don’t}
    pronounced \textit{doen’t}, not \textit{dewn’t}; \bad{amn’t}.
  \end{itemize}
\end{itemize}

\myslide{The Contraction Lexical Rule}
\begin{center}\footnotesize
  \avm{[\type*{pi-rule}
    input & < \2, [ syn & [ head & [\type*{verb} form & fin \\ aux & $+$ \\ pol & $-$ ] ] \\
                    arg-st & \tag{B} \\
                    sem & [ index & s$_1$ \\ restr & \tag{A} ] ] > \\
    output & < F$_{NEG}$(\2), [ syn & [ head & [ pol & $+$ ] \\ val & [ spr & < X > ] ] \\
                    arg-st & \tag{B} \\
                    sem & [ index & s$_2$ \\
                            restr & < [ reln & \textbf{not} \\ sit & s$_2$ \\ arg & s$_1$ ] > $\oplus$ \tag{A} ] ] > ]}
\end{center}
\emp{F$_{NEG}$ does the morphology: \lex{would} \into\ \lex{wouldn’t},
  \lex{will} \into\ \lex{won’t}.}

\qaslide{Most of the Work Is in the Semantics}{
\begin{center}\scriptsize
  \avm{[\type*{pi-rule}
    input & < \2, [ syn & [ head & [\type*{verb} form & fin \\ aux & $+$ \\ pol & $-$ ] ] \\
                    arg-st & \tag{B} \\
                    sem & [ index & \OV{s$_1$} \\ restr & \OV{\tag{A}} ] ] > \\
    output & < F$_{NEG}$(\2), [ syn & [ head & [ pol & $+$ ] \\ val & [ spr & < X > ] ] \\
                    arg-st & \tag{B} \\
                    sem & [ index & \OVB{s$_2$} \\
                            restr & < \OVB{[ reln & \textbf{not} \\ sit & s$_2$ \\ arg & s$_1$ ]} > $\oplus$ \OV{\tag{A}} ] ] > ]}
\end{center}
\begin{itemize}\addtolength{\itemsep}{-1ex}
\item Where does the \textbf{not} predication come from here?
  \answer{\\ {\small\emp{No separate word carries it, so the rule adds
    it to the verb's own \ft{restr} and makes $s_2$ the new \ft{index}:
    \eng{would not} packed into one word.}}}
\end{itemize}
}

\qaslide{What Does \ft{pol} Do?}{
\begin{center}\scriptsize
  \avm{[\type*{pi-rule}
    input & < \2, [ syn & [ head & [\type*{verb} form & fin \\ aux & $+$ \\ pol & \OV{$-$} ] ] \\
                    arg-st & \tag{B} \\
                    sem & [ index & s$_1$ \\ restr & \tag{A} ] ] > \\
    output & < F$_{NEG}$(\2), [ syn & [ head & [ pol & \OV{$+$} ] \\ val & [ spr & < X > ] ] \\
                    arg-st & \tag{B} \\
                    sem & [ \ldots ] ] > ]}
\end{center}
\begin{itemize}\addtolength{\itemsep}{-1ex}
\item Both the Contraction LR and the ADV$_{pol}$-Addition LR go from
  [\ft{pol} $-$] to [\ft{pol} $+$].  What does that buy us?
  \answer{\\ {\small\emp{Each rule can apply only once, and not after
    the other: it blocks \bad{We can’tn’t stop} and \bad{They won’t TOO
      mind}.  Verbs are [\ft{pol} $-$] unless a rule has polarized them.}}}
\end{itemize}
}

\myslide{Contraction: Sample Tree}
\begin{center}
\begin{forest}
cat
[{S} [{NP} [{Leslie}]]
  [{VP} [{V} [{wouldn’t}]] [{VP} [{eat the entire pizza}, roof]]]]
\end{forest}
\end{center}
\begin{itemize}\addtolength{\itemsep}{-1ex}
\item Exactly the tree for \eng{Leslie would eat the entire pizza}:
  the negation is inside the word \lex{wouldn’t}.
\end{itemize}

\qaslide{Ellipsis}{
\begin{itemize}
\item Ellipsis allows VPs to be omitted, so long as they would have
  been preceded by an auxiliary:
  \begin{exe}
  \ex \eng{Pat couldn’t have been watching us, but Chris could have been.}
  \ex \eng{Pat couldn’t have been watching us, but Chris could have.}
  \ex \eng{Pat couldn’t have been watching us, but Chris could.}
  \ex \bad{Pat couldn’t have been watching us, but Chris.}
  \end{exe}
\item Unlike the other NICE properties, this holds of all auxiliaries,
  not just finite ones.
\item What is the elliptical counterpart of a sentence with no
  auxiliary?  \eng{Whenever Pat watches TV, Chris watches TV.}
  \answer{\\ {\small\emp{\eng{Whenever Pat watches TV, Chris does.}
    Once more, \lex{do}.}}}
\end{itemize}
}

\qaslide{The Ellipsis Lexical Rule}{
\begin{center}
  \avm{[\type*{d-rule}
    input & < \1, [\type*{auxv-lxm} arg-st & < \2 > $\oplus$ \tag{A} ] > \\
    output & < \1, [\type*{dervv-lxm} arg-st & < \2 > ] > ]}
\end{center}
\begin{itemize}\addtolength{\itemsep}{-1ex}
\item Note that this is a derivational LR (\val{d-rule}), that is,
  lexeme-to-lexeme.
\item This means that \ft{syn} and \ft{sem} are unchanged, by default.
\item What is \val{dervv-lxm} and why do we have it?
  \answer{\\ {\small\emp{A leaf type of \val{verb-lxm} (derived verb
    lexeme) for the output: it is no longer an \val{auxv-lxm}, since it
    has no VP complement, and it must say [\ft{aux} $+$] so as not to
    inherit the default [\ft{aux} $-$] from \val{verb-lxm} (SWB
    p.\,420).}}}
\end{itemize}
}

\myslide{Ellipsis: Sample Output}
\begin{center}\small
  \avm{< \textnormal{could}, [\type*{dervv-lxm}
    syn & [ head & [ form & fin \\ aux & $+$ \\ pol & $-$ \\ agr & \1 ] \\
            val & [ spr & < [ agr & \1 ] > ] ] \\
    arg-st & < NP > \\
    sem & [ mode & prop \\ index & s$_1$ \\
            restr & < [ reln & \textbf{could} \\ sit & s$_1$ \\ arg & s$_2$ ] > ] ] >}
\end{center}
\emp{A finite \lex{could} with only a subject, and a meaning whose
  argument $s_2$ is not supplied by anything in the sentence.}

\myslide{Ellipsis: Sample Tree}
\begin{center}
\scalebox{0.9}{\begin{forest}
cat
[{S} [{NP} [{Kim}]]
  [{VP} [{V} [{could}]]
    [{VP} [{V} [{have}]]
      [{VP} [{V} [{been}]] [{VP} [{attending the conference}, roof]]]]]]
\end{forest}}
\end{center}
\begin{itemize}\addtolength{\itemsep}{-1ex}
\item The full sentence: three auxiliaries, each taking the next VP as
  its complement.
\end{itemize}

\myslide{Ellipsis: Sample Trees, with the Ellipsis LR}
\begin{center}
\begin{minipage}[t]{0.36\textwidth}\vspace{0pt}\centering
\scalebox{0.75}{\begin{forest}
cat
[{S} [{NP} [{Kim}]]
  [{VP} [{V} [{could}]] [{VP} [{V} [{have}]] [{VP} [{V} [{been}]]]]]]
\end{forest}}
\end{minipage}
\begin{minipage}[t]{0.3\textwidth}\vspace{0pt}\centering
\scalebox{0.75}{\begin{forest}
cat
[{S} [{NP} [{Kim}]]
  [{VP} [{V} [{could}]] [{VP} [{V} [{have}]]]]]
\end{forest}}
\end{minipage}
\begin{minipage}[t]{0.25\textwidth}\vspace{0pt}\centering
\scalebox{0.75}{\begin{forest}
cat
[{S} [{NP} [{Kim}]] [{VP} [{V} [{could}]]]]
\end{forest}}
\end{minipage}
\end{center}
\begin{itemize}\addtolength{\itemsep}{-1ex}
\item The last auxiliary in each tree is a \val{dervv-lxm}: its VP
  argument is gone, so the Head-Complement Rule applies with no
  complements.
\end{itemize}

\qaslide{Semantics of Ellipsis}{
\begin{center}
\scalebox{0.7}{\begin{forest}
cat
[{S} [{NP} [{Kim}]] [{VP} [{V} [{could}]]]]
\end{forest}}
\end{center}
\vspace{-2ex}
\begin{itemize}\addtolength{\itemsep}{-1ex}
\item What is the \ft{sem} value of the S node of this tree?
  \answer{\\
  \begin{center}\footnotesize
    \avm{[ index & s$_1$ \\ mode & prop \\
           restr & < [ reln & \textbf{name} \\ name & Kim \\ named & i ],
                     [ reln & \textbf{could} \\ sit & s$_1$ \\ arg & s$_2$ ] > ]}
  \end{center}
  {\small\emp{Note: $s_2$ has to be filled in by context, which is
    exactly what ellipsis feels like.}}}
\end{itemize}
}

\myslide{Infinitival \lex{to} Revisited}
\begin{itemize}
\item VP ellipsis can occur after \lex{to}:
  \begin{exe}
  \ex \eng{We didn’t find the solution, but we tried to.}
  \end{exe}
\item This is covered by our Ellipsis LR if we say \lex{to} is
  [\ft{aux} $+$].
\item Since \ft{aux} is declared on type \val{verb}, it follows that
  \lex{to} is a verb.
\item That was the promise from Lecture 12: now kept.
\end{itemize}

\myslide{\lex{do} Revisited}
\begin{itemize}\addtolength{\itemsep}{-0.5ex}
\item Chomsky’s old analysis: in sentences without auxiliaries
  \begin{itemize}
  \item Tense can get separated from the verb in various ways:
    negation/reaffirmation inserts something between Tense and the
    following verb; inversion moves Tense to the left of the subject NP;
    ellipsis deletes what follows Tense.
  \item When this happens, \lex{do} is inserted to support Tense.
  \end{itemize}
\item Our counterpart
  \begin{itemize}
  \item NICE properties hold only of auxiliaries.
  \item \lex{do} is a semantically empty auxiliary.
  \item So, negated, reaffirmed, inverted and elliptical sentences that
    are the semantic counterparts to sentences without auxiliaries are
    ones with \lex{do}.
  \end{itemize}
\end{itemize}

\myslide{Summary}
\begin{itemize}
\item Our analysis employs straightforward mechanisms
  \begin{itemize}
  \item lexical entries for auxiliaries (a subtype of raising verbs)
  \item 3 new features (\ft{aux}, \ft{pol}, \ft{inv})
  \item 4 lexical rules (ADV$_{pol}$-Addition, Inversion, Contraction,
    Ellipsis)
  \end{itemize}
\item We handle a complex array of facts
  \begin{itemize}
  \item co-occurrence restrictions (ordering and iteration)
  \item the NICE properties
  \item auxiliary \lex{do}
  \item combinations of NICE constructions
  \end{itemize}
\end{itemize}

\section[Step by step: \eng{Kim did not leave} and \eng{Did Kim leave?}]{Step by Step}

\myslide{Step by Step: What \lex{did} Asks For}
\begin{center}
\scalebox{0.7}{\begin{forest}
cat, for tree={align=center}
[{S}
[{\avmbox{1}NP} [{Kim}]]
[{VP\\ {\small$[$ \ft{spr} $\langle \avmbox{1} \rangle$, \ft{comps} $\langle\ \rangle$ $]$}}
[{V\\ {\small$[$ \ft{spr} $\langle \avmbox{1} \rangle$,}\\ {\small\ft{comps} $\langle \avmbox{2}, \avmbox{3} \rangle$ $]$}} [{did}]]
[{\avmbox{2}ADV$_{pol}$} [{not}]]
[{\avmbox{3}VP\\ {\small$[$ \ft{form} base, \ft{aux} $-$, \ft{spr} $\langle \avmbox{1} \rangle$ $]$}} [{leave}]]]]
\end{forest}}
\end{center}
\vspace{-2ex}
\begin{itemize}\addtolength{\itemsep}{-1.5ex}
\item The lexeme \lex{do} has \ft{arg-st} $\langle$ \avmbox{1},
  \avmbox{3}VP[\ft{form} base, \ft{aux} $-$, \ft{spr} $\langle$
  \avmbox{1} $\rangle$] $\rangle$: a raising verb.
\item The ADV$_{pol}$-Addition LR inserts \avmbox{2}: $\langle$
  \avmbox{1}, \avmbox{2}, \avmbox{3} $\rangle$; the ARP puts \avmbox{1}
  on \ft{spr} and the rest on \ft{comps}.
\item \lex{leave} still wants \avmbox{1} as its subject: Kim is the
  leaver, \lex{did} adds nothing, \lex{not} negates.
\end{itemize}

\qaslide{Which Lexical Rules License Each Word?}{
\eng{Kim did not leave.}
\begin{itemize}\addtolength{\itemsep}{-1ex}
\item \eng{Kim}?
  \answer{\\ {\small\emp{Constant Lexeme LR (\val{pn-lxm}).}}}
\item \eng{did}?
  \answer{\\ {\small\emp{Two rules: the Past-Tense Verb LR (lexeme
    \into\ word, \lex{do} \into\ \lex{did}, [\ft{form} fin] as the
    lexeme already requires), then the ADV$_{pol}$-Addition LR (word
    \into\ word, a \val{pi-rule}), which adds \lex{not} to \ft{arg-st}
    and makes it [\ft{pol} $+$].}}}
\item \eng{not}?
  \answer{\\ {\small\emp{Constant Lexeme LR, from the ADV$_{pol}$ entry.}}}
\item \eng{leave}?
  \answer{\\ {\small\emp{Base Form LR: \lex{did} asks for [\ft{form}
    base].}}}
\end{itemize}
}

\myslide{Step by Step: \eng{Did Kim leave?}}
\begin{center}
\scalebox{0.75}{\begin{forest}
cat, for tree={align=center}
[{S\\ {\small$[$ \ft{spr} $\langle\ \rangle$, \ft{comps} $\langle\ \rangle$, \ft{mode} ques $]$}}
[{V\\ {\small$[$ \ft{inv} $+$, \ft{spr} $\langle\ \rangle$,}\\ {\small\ft{comps} $\langle \avmbox{1}, \avmbox{2} \rangle$ $]$}} [{Did}]]
[{\avmbox{1}NP} [{Kim}]]
[{\avmbox{2}VP\\ {\small$[$ \ft{form} base, \ft{spr} $\langle \avmbox{1} \rangle$ $]$}} [{leave}]]]
\end{forest}}
\end{center}
\vspace{-2ex}
\begin{itemize}\addtolength{\itemsep}{-1.5ex}
\item Same \ft{arg-st} as before the rule: $\langle$ \avmbox{1},
  \avmbox{2} $\rangle$.  The Inversion LR only sets \ft{spr} to
  $\langle\ \rangle$ and \ft{mode} to \val{ques}.
\item So the ARP puts both arguments on \ft{comps}, and the subject
  follows the verb.
\item \avmbox{2} still has \ft{spr} $\langle$ \avmbox{1} $\rangle$:
  Kim is the leaver, agreement and case are checked as usual.
\end{itemize}

\qaslide{Which Rule Licenses Each Node?}{
\begin{center}
\scalebox{0.6}{\begin{forest}
cat
[{S} [{NP} [{Kim}]]
  [{VP} [{V} [{did}]] [{ADV$_{pol}$} [{not}]] [{VP} [{V} [{leave}]]]]]
\end{forest}}
\hspace{2em}
\scalebox{0.6}{\begin{forest}
cat
[{S} [{V} [{Did}]] [{NP} [{Kim}]] [{VP} [{V} [{leave}]]]]
\end{forest}}
\end{center}
\vspace{-2ex}
\begin{itemize}\addtolength{\itemsep}{-1.5ex}
\item Top S of \eng{Kim did not leave}?
  \answer{{\small\emp{Head-Specifier Rule: the VP and its specifier \eng{Kim}.}}}
\item VP \eng{did not leave}?
  \answer{{\small\emp{Head-Complement Rule: \eng{did} and both complements.}}}
\item VP \eng{leave}?
  \answer{{\small\emp{Head-Complement Rule with no complements.}}}
\item S of \eng{Did Kim leave?}
  \answer{{\small\emp{Head-Complement Rule: \eng{Did} and both
    complements; \ft{spr} and \ft{comps} are empty, so it is an S.}}}
\end{itemize}
}

\myslide{If You Want to Try}
\begin{itemize}
\item There is no problem set for this lecture.
\item If you would like to practise, \citet[Chapter 13]{sag:wasow:bender:2003}
  Problems 1--3 cover the ground above:
  \begin{itemize}
  \item Problem 1: why there are no modal imperatives.
  \item Problem 2: \lex{do} support without a \lex{do}-insertion rule.
  \item Problem 3: trees for inverted sentences.
  \end{itemize}
\item Try \eng{You could do that, but you better not}: draw the tree,
  and say which lexical type and which lexical rules license each word.
\end{itemize}

\overview
\MyLogo{\citet{sag:wasow:bender:2003} --- Chapter 13}

\myslide{Acknowledgments and References}

\begin{itemize}
\item Course design and slides borrow heavily from Emily Bender's course:
\textit{Linguistics 566: Introduction to Syntax for Computational Linguistics}
\\ \url{http://courses.washington.edu/ling566}
\item This lecture complements \citet[Chapter 13]{sag:wasow:bender:2003}
\end{itemize}
\small
\bibliographystyle{aclnat}
\bibliography{abb,mtg,nlp,ling}
\end{document}

%%% Local Variables:
%%% coding: utf-8
%%% mode: latex
%%% TeX-PDF-mode: t
%%% TeX-engine: xetex
%%% End:
